% META: {"id":"TEXP-MAT-CO-002","chapitre":"TEXP-MATRICES-MARKOV","type_objet":"corrige","exercice_ref":"TEXP-MAT-EX-002","capacites_codes":["C2"],"sources_inspiration":[],"status":"approved","origin":"generated","status_history":["generated","audited","accepted","approved"],"release_acceptance":"RELEASE_OWNER_BATCH_ACCEPTANCE","acceptance_closure_digest":"sha256:9b3ccf9a81c5520fbb7e03b7b2d3e7bf2057a02834b6d908bf3be5f49f3b553f","acceptance_receipt_digest":"sha256:4fad86f0282e0fa4e0419cca1ad6379ae7af5a280c04a4a90880f90a1c044242"}
% BEGIN-VERIFY
% from sympy import Matrix
% qte = Matrix([[10,5,20]])
% prix = Matrix([[12],[15],[8]])
% total = qte*prix
% assert total == Matrix([[355]])
% END-VERIFY
\begin{corrige}{TEXP-MAT-EX-002}
\textbf{1.} $Q=\begin{pmatrix}10&5&20\end{pmatrix}$ (quantités de romans, BD, mangas) et $T=\begin{pmatrix}12\\15\\8\end{pmatrix}$ (prix unitaires correspondants).

\textbf{2.} $Q\times T=10\times12+5\times15+20\times8=120+75+160=355$. Ce nombre représente la \textbf{recette totale} (en euros) de la journée.
\end{corrige}
